Let $G$ be a graph and $H$ a minor of $G$. Every vertex of $G_i$ is in a part of $\mathcal{P}_i$.
The edges in black are from the original graph $G$, the edges in red are
from the flip $G_i$. The partition is $\mathcal{P}_i$, and the set of its
parts is $\mathcal{Q}$. Here $Q_u = Q_1$. Every part $P_u$ is flipped with
$P_W$. The bound is $d^2 t^3$, and it holds for each $r$. So the class is
$\mathcal{C}$, its closure is $\overline{\mathcal{C}}$, and the result is
$x_{i+1}$. Each of $u$, $v$ and $w$ lies in $X$. The radius is $r$.
